\documentclass[12pt,a4paper,oneside]{report}

\usepackage[utf8]{inputenc}
\usepackage[T1]{fontenc}
\usepackage[spanish,es-tabla]{babel}
\usepackage{csquotes}
\usepackage{geometry}
\usepackage{graphicx}
\usepackage{booktabs}
\usepackage{longtable}
\usepackage{array}
\usepackage{hyperref}
\usepackage{xcolor}
\usepackage{setspace}
\usepackage{titlesec}
\usepackage[backend=biber,style=apa,sorting=nyt]{biblatex}

\geometry{
    left=3cm,
    right=2.5cm,
    top=3cm,
    bottom=2.5cm
}

\onehalfspacing
\graphicspath{{figures/}}
\addbibresource{bibliografia.bib}

\hypersetup{
    colorlinks=true,
    linkcolor=black,
    citecolor=black,
    urlcolor=blue,
    pdftitle={Informe de Tesina},
    pdfauthor={Evelin Emilce Aragon}
}

\titleformat{\chapter}
  {\normalfont\Large\bfseries}
  {\thechapter.}
  {0.7em}
  {}

\titleformat{\section}
  {\normalfont\large\bfseries}
  {\thesection}
  {0.7em}
  {}


\title{Traductor móvil en tiempo real para Lengua de Señas Argentina basado en aprendizaje automático y ejecución local en Android}
\author{Evelin Emilce Aragón}
\date{\today}

\begin{document}

\begin{titlepage}
    \centering
    {\Large Universidad Nacional de General Sarmiento\par}
    \vspace{1.5cm}
    {\large Licenciatura en Sistemas\par}
    \vspace{2.5cm}
    {\Huge\bfseries Traductor móvil en tiempo real para Lengua de Señas Argentina\par}
    \vspace{0.4cm}
    {\Large basado en aprendizaje automático y ejecución local en Android\par}
    \vspace{2.5cm}
    {\large Autora: Evelin Emilce Aragón\par}
    \vspace{0.3cm}
    \vfill
    {\large \today\par}
\end{titlepage}

\pagenumbering{roman}

\chapter*{Resumen}
\addcontentsline{toc}{chapter}{Resumen}
Este trabajo presenta SeÑAR, un prototipo de aplicación Android que facilita la comunicación entre personas sordas y oyentes: reconoce en tiempo real doce señas de Lengua de Señas Argentina (LSA) y las traduce a texto y voz, y transcribe a texto la voz de la persona oyente. Todo el procesamiento ocurre en el dispositivo, sin conexión a internet.

La aplicación detecta las manos con MediaPipe, construye una secuencia de 15 cuadros con las coordenadas de 21 puntos por mano y la clasifica con una red recurrente de 236\,KB ejecutada con TensorFlow Lite. El conjunto de datos reúne tomas de cámara web y 449 tomas de tres personas grabadas con la propia aplicación, siguiendo un protocolo que incluye cuatro condiciones adversas. Se incorporó además un esquema de aprendizaje federado en el que los dispositivos entrenan la parte final del modelo con sus propias tomas y envían solo una actualización de unos 13\,KB, que un servidor de la red local combina y redistribuye.

Evaluado con personas que no vio durante el entrenamiento, el modelo acierta el 67,3\,\% de las señas, frente al 94,3\,\% que se obtiene mezclando las tomas al azar: la diferencia muestra que la diversidad de personas pesa más que la cantidad de repeticiones. Con la confianza calibrada, la aplicación acierta el 83\,\% de lo que muestra. De las condiciones adversas, solo el cuerpo girado produce una caída significativa, atribuible a la detección de las manos. La aplicación funciona en seis dispositivos, entre 6,6 y 24,9 cuadros por segundo, incluido un teléfono de gama baja de 2019, y una prueba en modo avión confirmó que la traducción de señas y la transcripción de voz funcionan sin conexión. El piloto federado con dos dispositivos reales mostró que el ciclo completo es viable y que la personalización local complementa al modelo compartido.

\revisar{¿dejar o sacar las palabras clave? Las tesinas de referencia y la plantilla UNGS no las incluyen.}{\textbf{Palabras clave:} Lengua de Señas Argentina, accesibilidad, visión por computadora, aprendizaje automático, Android, TensorFlow Lite.}

\chapter*{Agradecimientos}
\addcontentsline{toc}{chapter}{Agradecimientos}
Completar al finalizar el informe.

\chapter*{Lista de abreviaciones}
\addcontentsline{toc}{chapter}{Lista de abreviaciones}

% Mantener en orden alfabético. Agregar cada sigla la primera vez que se use en el texto.
\begin{description}
    \item[CPU] \emph{Central Processing Unit} (unidad central de procesamiento).
    \item[FedAvg] \emph{Federated Averaging} (promedio federado).
    \item[FPS] \emph{Frames per second} (cuadros por segundo).
    \item[IID] Independientes e idénticamente distribuidos (referido a datos).
    \item[LSA] Lengua de Señas Argentina.
    \item[LSTM] \emph{Long Short-Term Memory} (memoria de corto y largo plazo).
    \item[RGB] \emph{Red, Green, Blue} (modelo de color rojo, verde y azul de una imagen).
    \item[RNN] \emph{Recurrent Neural Network} (red neuronal recurrente).
    \item[TFLite] TensorFlow Lite.
    \item[TTS] \emph{Text-to-Speech} (síntesis de voz a partir de texto).
\end{description}

\chapter*{Glosario}
\addcontentsline{toc}{chapter}{Glosario}

% Mantener en orden alfabético. Agregar cada término a medida que aparezca en los capítulos.
\begin{description}
    \item[Aprendizaje federado] Técnica de aprendizaje automático en la que varios dispositivos entrenan un modelo compartido a partir de sus datos locales y envían solo las actualizaciones del modelo, sin transferir los datos crudos.
    \item[Cuantización] Técnica que reduce la precisión numérica de los parámetros de un modelo (por ejemplo, de punto flotante de 32 bits a enteros de 8 bits) para disminuir su tamaño y acelerar la inferencia.
    \item[Estimación de pose] Técnica de visión por computadora que determina la posición de un conjunto de puntos clave del cuerpo o de las manos a partir de una imagen.
    \item[Glosa] Transcripción de una seña mediante una palabra de una lengua oral, utilizada como representación intermedia en la traducción de lenguas de señas.
    \item[Handedness] Lateralidad de una mano detectada, es decir, si se trata de la mano izquierda o de la derecha.
    \item[Inferencia] Uso de un modelo ya entrenado para obtener una predicción a partir de una entrada nueva.
    \item[Landmark] Punto de referencia que representa la posición de una articulación o extremo de la mano dentro de la imagen.
    \item[Pipeline] Secuencia de etapas de procesamiento en la que la salida de cada etapa es la entrada de la siguiente.
    \item[Seña aislada] Seña individual que se reconoce de forma independiente, sin considerar las señas anteriores o posteriores dentro de una oración.
\end{description}


\clearpage
\phantomsection
\addcontentsline{toc}{chapter}{Índice de figuras}
\listoffigures

\clearpage
\phantomsection
\addcontentsline{toc}{chapter}{Índice de cuadros}
\listoftables

\tableofcontents
\clearpage

\pagenumbering{arabic}

\chapter{Introduccion}

\section{Contexto y motivacion}

La comunicacion accesible constituye un aspecto central para la inclusion social y tecnologica. En el caso de las personas sordas, la Lengua de Señas Argentina (LSA) representa una forma de comunicacion natural y fundamental, pero su uso no siempre encuentra correspondencia en herramientas digitales ampliamente disponibles. Esta situacion genera barreras de interaccion en distintos contextos cotidianos, especialmente cuando no existe un interlocutor que conozca la lengua de señas.

En este marco, el presente trabajo aborda el desarrollo de un prototipo de software orientado a la traduccion de señas a texto y voz mediante un dispositivo movil Android. La propuesta se centra en una arquitectura de ejecucion local, sin dependencia permanente de conexion a internet, con el objetivo de favorecer la portabilidad, la privacidad y la disponibilidad del sistema en escenarios de uso real.

El proyecto combina tecnicas de vision por computadora y aprendizaje automatico. En particular, utiliza la deteccion de landmarks de manos para representar geometricamente los gestos capturados por la camara y un modelo TensorFlow Lite para su clasificacion en tiempo real. Sobre esa base, la aplicacion produce una salida textual y sonora que busca facilitar la comunicacion entre personas sordas y oyentes.

\section{Problema}

A pesar del crecimiento de las tecnologias de asistencia, siguen siendo limitadas las soluciones orientadas a la traduccion automatica de Lengua de Señas Argentina en dispositivos moviles, especialmente aquellas que puedan ejecutarse de forma local. Muchas propuestas dependen de servicios externos, conectividad permanente o conjuntos de datos que no representan adecuadamente las particularidades de la LSA. Esto reduce su aplicabilidad en contextos reales y plantea desafios vinculados con la latencia, la privacidad y la accesibilidad.

Frente a esta situacion, surge la necesidad de explorar alternativas que permitan reconocer un conjunto acotado de señas en tiempo real sobre dispositivos Android, aprovechando modelos livianos y tecnicas de procesamiento eficientes. El problema que orienta esta tesina puede expresarse de la siguiente manera: ¿es factible desarrollar un traductor movil de LSA a texto y voz que funcione localmente en Android y ofrezca un comportamiento aceptable en condiciones realistas de uso?

\section{Objetivo general}

Desarrollar y evaluar un prototipo de aplicacion movil para Android capaz de traducir señas de Lengua de Señas Argentina a texto y voz en tiempo real, mediante una arquitectura de ejecucion local basada en vision por computadora y aprendizaje automatico.

\section{Objetivos especificos}

\begin{itemize}
    \item Diseñar una arquitectura de procesamiento apta para ejecutarse en dispositivos moviles Android.
    \item Implementar la captura de imagen y el analisis de gestos a partir de la camara del dispositivo.
    \item Detectar landmarks de manos utilizando herramientas de vision por computadora.
    \item Construir una representacion de caracteristicas adecuada para la clasificacion temporal de señas.
    \item Integrar un modelo TensorFlow Lite para reconocer un conjunto inicial de palabras o expresiones.
    \item Generar una salida en texto y voz para mejorar la interpretacion del resultado.
    \item Evaluar el comportamiento del prototipo en distintos dispositivos y analizar sus limitaciones.
\end{itemize}

\section{Alcance}

El trabajo se focaliza en el desarrollo de un prototipo funcional orientado a validar la factibilidad tecnica de la propuesta. No se busca resolver la traduccion completa de la Lengua de Señas Argentina ni cubrir toda su riqueza linguistica, sino implementar un sistema capaz de reconocer un conjunto limitado de señas previamente definidas. Asimismo, el enfoque esta puesto en la traduccion de seña a texto y voz, quedando fuera del alcance la traduccion completa en sentido inverso o la interpretacion de estructuras gramaticales complejas.

\section{Metodologia de trabajo}

La metodologia adoptada fue de tipo aplicada y experimental. En una primera etapa se relevaron antecedentes sobre reconocimiento de señas, accesibilidad y ejecucion de modelos en dispositivos moviles. Luego se diseño la arquitectura general del sistema, definiendo el flujo de captura, deteccion, extraccion de caracteristicas, clasificacion e interfaz de usuario. Finalmente, se implemento el prototipo y se realizaron pruebas de funcionamiento y rendimiento en distintos dispositivos Android para analizar su viabilidad practica.

\section{Estructura del informe}

El presente informe se organiza de la siguiente manera. En el Capitulo 2 se desarrolla el marco teorico y se revisan antecedentes relacionados con la traduccion automatica de lenguas de señas. En el Capitulo 3 se delimita el problema abordado y se especifican los requisitos del sistema. En el Capitulo 4 se presenta la propuesta de solucion y la arquitectura general del prototipo. En el Capitulo 5 se describe la implementacion realizada. En el Capitulo 6 se exponen las pruebas efectuadas y los resultados obtenidos. Finalmente, en el Capitulo 7 se presentan las conclusiones, limitaciones y posibles lineas de trabajo futuro.

\chapter{Marco teorico y antecedentes}

\section{Accesibilidad y comunicacion en Lengua de Señas Argentina}

Completar con una descripcion de la LSA, su importancia como lengua natural de la comunidad sorda argentina y las barreras comunicacionales que pueden aparecer en la interaccion con personas oyentes.

\section{Sistemas de traduccion de señas}

Completar con antecedentes sobre sistemas de reconocimiento y traduccion de lenguas de señas, diferenciando enfoques basados en video, sensores, guantes, landmarks, glosas o representaciones intermedias.

\section{Vision por computadora para reconocimiento gestual}

Completar con los conceptos principales de vision por computadora aplicados al reconocimiento de gestos: captura de imagen, deteccion de manos, seguimiento temporal y representacion de movimiento.

\section{Landmarks de manos}

Completar con la idea de landmarks como puntos de referencia que representan la posicion de articulaciones de la mano. En este proyecto, estos puntos permiten transformar una imagen en una representacion numerica apta para un modelo de clasificacion.

\section{Modelos livianos en dispositivos moviles}

Completar con la importancia de ejecutar modelos livianos en dispositivos moviles, considerando restricciones de CPU, memoria, bateria, latencia y disponibilidad offline.

\section{Limitaciones de los enfoques existentes}

Completar con las principales limitaciones identificadas: escasez de datasets en LSA, dependencia de servidores externos, problemas de privacidad, necesidad de conectividad y dificultad de generalizar gestos entre distintas personas.

\chapter{Metodología}
\label{cap:metodologia}

\section{Tipo de trabajo}

\pendiente{Caracterizar el trabajo como desarrollo tecnológico con evaluación experimental: se construye un prototipo y se valida empíricamente su viabilidad en dispositivos móviles reales.}

\section{Etapas del trabajo}

\pendiente{Describir las etapas efectivamente seguidas (captura de datos, entrenamiento, integración en Android, pruebas), contrastándolas con el plan de trabajo de la propuesta aprobada.}

\section{Herramientas utilizadas}

\pendiente{Listar y justificar brevemente las herramientas de entrenamiento (Python, TensorFlow/Keras, MediaPipe) y de desarrollo móvil (Android, Kotlin, Jetpack Compose, CameraX, TensorFlow Lite, Text-to-Speech).}

\section{Métricas y criterios de evaluación}

\pendiente{Definir las métricas del modelo (accuracy, F1 macro y micro, matriz de confusión, consistencia estadística $\sigma$) y las métricas de rendimiento en el dispositivo (FPS efectivos, latencia de inferencia, tamaño del modelo).}

\section{Diseño de las pruebas}

\pendiente{Describir las condiciones de prueba: usuarios, iluminación, ángulo y distancia de cámara, y dispositivos utilizados.}

\chapter{Construcción del conjunto de datos y del modelo}
\label{cap:dataset-modelo}

\section{Selección del vocabulario}

\pendiente{Justificar las señas elegidas y su cantidad. Validar contra el vocabulario vigente al cerrar el entrenamiento.}

\section{Protocolo de captura}

\pendiente{Describir cómo se capturaron las muestras: personas, repeticiones, condiciones de luz, posición de cámara, herramienta de captura y cantidad total de muestras.}

\section{Representación de las señas}

\pendiente{Explicar la representación de cada seña como secuencia de $T$ cuadros con 126 valores por cuadro (21 puntos por mano, dos manos, tres coordenadas), es decir, una matriz $\mathbf{X} \in \mathbb{R}^{T \times 126}$.}

\section{Preprocesamiento y normalización}

\pendiente{Describir la normalización temporal a longitud fija, el tratamiento de manos ausentes y cualquier normalización espacial aplicada.}

\section{Arquitectura del modelo}

\pendiente{Describir la arquitectura final y justificarla. Contrastar con la arquitectura LSTM planteada en la propuesta si hubo cambios.}

\section{Entrenamiento}

\pendiente{Partición de datos, hiperparámetros, regularización y curvas de entrenamiento.}

\section{Conversión a TensorFlow Lite y cuantización}

\pendiente{Describir la conversión del modelo y comparar variantes (por ejemplo, float32 frente a optimizada) en tamaño y desempeño.}

\chapter{Diseño e implementación de la aplicación}
\label{cap:diseno}

Este capítulo describe la aplicación SeÑAR tal como quedó al cierre del trabajo: los requisitos que guiaron su diseño, la arquitectura, el recorrido de cada cuadro de la cámara hasta la palabra que se muestra y se pronuncia, y los modos complementarios de voz a texto y de entrenamiento. Los problemas que obligaron a cambiar el diseño en el camino se tratan en el Capítulo~\ref{cap:desafios}, y las mediciones de rendimiento, en el Capítulo~\ref{cap:resultados}.

\section{Requisitos del sistema}

\subsection{Requisitos funcionales}

Los requisitos comprometidos en la propuesta son los siguientes:

\begin{itemize}
    \item Capturar imagen desde la cámara del dispositivo móvil.
    \item Detectar las manos y sus puntos de referencia (\emph{landmarks}) en tiempo real.
    \item Construir una secuencia temporal de características.
    \item Clasificar la seña realizada.
    \item Mostrar el resultado en pantalla.
    \item Reproducir el resultado mediante síntesis de voz.
    \item Transcribir voz a texto para la persona oyente.
\end{itemize}

Durante el desarrollo se sumaron tres funciones que no estaban en la propuesta y que surgieron de la necesidad de obtener datos reales: un modo de entrenamiento para grabar tomas etiquetadas desde la propia aplicación, una sesión guiada que ordena la grabación según un protocolo, y el aprendizaje federado (Capítulo~\ref{cap:federado}).

\subsection{Requisitos no funcionales}

\begin{itemize}
    \item Ejecutar el reconocimiento de señas y la transcripción de voz en el dispositivo, sin conexión a internet.
    \item No enviar imágenes ni audio a servidores externos.
    \item Mantener tiempos de respuesta aceptables para una conversación: la captura de una seña no debería superar alrededor de un segundo y medio en un teléfono de gama media.
    \item Funcionar en teléfonos de gama media y baja, no solo en equipos de gama alta.
    \item Ser utilizable por personas sin conocimientos técnicos.
\end{itemize}

\section{Arquitectura general}

La aplicación está escrita en Kotlin para Android y organizada en paquetes con una responsabilidad cada uno (Tabla~\ref{tab:paquetes}). El reconocimiento de señas se estructura como un pipeline que se ejecuta íntegramente en el dispositivo: la cámara entrega cuadros, MediaPipe detecta las manos, la aplicación construye una secuencia de características y un modelo de TensorFlow Lite la clasifica.

\begin{table}[h]
    \centering
    \small
    \begin{tabular}{lp{10cm}}
        \toprule
        Paquete & Responsabilidad \\
        \midrule
        \texttt{camera} & Configuración de CameraX y conversión de cada cuadro a una imagen derecha. \\
        \texttt{detection} & Detector de manos de MediaPipe, con elección de GPU o CPU. \\
        \texttt{features} & Vector de 126 valores por cuadro, secuencia temporal e implementación en Kotlin del contrato de características. \\
        \texttt{ml} & Máquina de estados de la captura, clasificador, etiquetas y calibración. \\
        \texttt{speech} & Síntesis de voz con el motor de Android. \\
        \texttt{voice} & Reconocimiento de voz y limpieza del texto transcripto. \\
        \texttt{dataset} & Formato de las tomas de entrenamiento y seudónimo de quien graba. \\
        \texttt{federado} & Paquetes de modelo, entrenamiento local y comunicación con el servidor. \\
        \texttt{settings} & Preferencias de la persona usuaria. \\
        \texttt{ui} & Interfaz en Jetpack Compose: navegación, pantallas y sistema de diseño. \\
        \bottomrule
    \end{tabular}
    \caption{Paquetes de la aplicación.}
    \label{tab:paquetes}
\end{table}

\section{Flujo de procesamiento}
\label{sec:flujo}

La Figura~\ref{fig:flujo} resume el recorrido de un cuadro de la cámara hasta la palabra reconocida.

\begin{figure}[h]
    \centering
    \begin{tikzpicture}[
        caja/.style={draw, rounded corners, align=center, minimum height=1.1cm, text width=2.6cm, font=\small},
        flecha/.style={-{Stealth}, thick}
    ]
        \node[caja] (cam) {Cámara\\(CameraX, RGBA)};
        \node[caja, right=0.7cm of cam] (rot) {Imagen derecha\\(rotación en la copia)};
        \node[caja, right=0.7cm of rot] (mp) {MediaPipe\\Hand Landmarker};
        \node[caja, below=1.0cm of mp] (vec) {Vector de 126\\valores por cuadro};
        \node[caja, left=0.7cm of vec] (seq) {Secuencia de\\10, 15 o 20 cuadros};
        \node[caja, left=0.7cm of seq] (tfl) {Modelo TFLite\\$15 \times 126 \rightarrow 13$};
        \node[caja, below=1.0cm of tfl] (cal) {Calibración\\y umbral};
        \node[caja, right=0.7cm of cal] (txt) {Texto\\en pantalla};
        \node[caja, right=0.7cm of txt] (voz) {Voz\\(TTS)};
        \draw[flecha] (cam) -- (rot);
        \draw[flecha] (rot) -- (mp);
        \draw[flecha] (mp) -- (vec);
        \draw[flecha] (vec) -- (seq);
        \draw[flecha] (seq) -- (tfl);
        \draw[flecha] (tfl) -- (cal);
        \draw[flecha] (cal) -- (txt);
        \draw[flecha] (txt) -- (voz);
    \end{tikzpicture}
    \caption{Recorrido de un cuadro en el modo de traducción de señas.}
    \label{fig:flujo}
\end{figure}

El pipeline no clasifica de forma continua, sino seña por seña. Una máquina de estados decide cuándo empieza y termina cada captura:

\begin{enumerate}
    \item \textbf{En espera.} No hay manos en el cuadro.
    \item \textbf{Preparación.} Aparecen manos y se mantienen durante 500\,ms. La espera evita capturar un movimiento casual al entrar al encuadre.
    \item \textbf{Captura.} Se acumulan cuadros con manos hasta completar la cantidad prevista. Si las manos desaparecen, la captura se descarta y se vuelve a la espera.
    \item \textbf{Clasificación.} La secuencia se lleva a 15 cuadros y se clasifica. Para habilitar la siguiente seña, la persona tiene que bajar las manos durante 600\,ms.
\end{enumerate}

La cantidad de cuadros por captura se ajusta a la velocidad del dispositivo: se capturan 10 cuadros si el dispositivo procesa menos de 8 por segundo, 20 si procesa más de 20 y 15 en los demás casos. De este modo, una seña dura aproximadamente lo mismo en un teléfono lento que en uno rápido, y el modelo siempre recibe la misma longitud.

\section{Decisiones tecnológicas}

\begin{itemize}
    \item \textbf{Android y Kotlin.} El público al que apunta la aplicación usa mayormente teléfonos Android, y Kotlin es el lenguaje recomendado para la plataforma. La interfaz se escribió con Jetpack Compose, el sistema declarativo actual de Android.
    \item \textbf{CameraX.} Ofrece un caso de uso de análisis de imágenes que entrega los cuadros a un hilo propio y resuelve las diferencias entre fabricantes, que son frecuentes en el acceso directo a la cámara.
    \item \textbf{MediaPipe Hand Landmarker.} Detecta hasta dos manos y entrega 21 puntos tridimensionales por mano, con un modelo de unos 8\,MB que funciona en el dispositivo y puede usar la GPU~\cite{zhang2020mediapipehands}. Es además el mismo detector que se usa en Python para extraer los landmarks de entrenamiento, lo que evita diferencias entre ambos extractores.
    \item \textbf{TensorFlow Lite.} Ejecuta el modelo entrenado en Keras sin conexión y con un intérprete pequeño. El modelo se exportó solo con operaciones básicas para no depender de la biblioteca de operaciones adicionales (Capítulo~\ref{cap:desafios}).
    \item \textbf{Servicios de voz del sistema.} La síntesis (\emph{TextToSpeech}) y el reconocimiento de voz (\emph{SpeechRecognizer}) son los de Android, que pueden funcionar sin conexión una vez descargado el paquete del idioma, y no agregan peso a la aplicación.
\end{itemize}

\subsection{Justificación de la ejecución local}

La ejecución local permite reducir la dependencia de conectividad, mejorar la disponibilidad de la aplicación en diferentes contextos y evitar el envío de imágenes o datos sensibles a servidores externos. Esta decisión resulta especialmente relevante en aplicaciones vinculadas con accesibilidad y comunicación personal.

La traducción de señas y la transcripción de voz no usan la red. La aplicación tiene, sin embargo, dos funciones que sí se comunican, ambas opcionales y fuera del uso de traducción:

\begin{itemize}
    \item \textbf{Cuenta de usuario.} El modo de entrenamiento requiere iniciar sesión, mediante el servicio de autenticación Clerk, para asociar las tomas a un seudónimo estable de la persona. La traducción y el diccionario se pueden usar sin cuenta.
    \item \textbf{Aprendizaje federado.} La pantalla correspondiente se comunica con un servidor de la red local, y solo cuando la persona lo solicita explícitamente (Capítulo~\ref{cap:federado}).
\end{itemize}

\section{Implementación del pipeline de reconocimiento}

\subsection{Captura de imagen con CameraX}

La cámara se configura con el caso de uso de análisis de imágenes y la estrategia de conservar solo el cuadro más reciente: si el procesamiento de un cuadro todavía no terminó cuando llega el siguiente, el anterior se descarta en lugar de encolarse. Así, en un dispositivo lento la aplicación procesa menos cuadros por segundo, pero nunca acumula retraso respecto de lo que ocurre frente a la cámara.

La resolución de análisis es configurable en tres niveles (360$\times$480, 480$\times$640 y 720$\times$960); el valor por defecto es el intermedio. La cámara entrega cada cuadro en formato RGBA y la aplicación lo copia a una imagen ya rotada a la orientación vertical, en una sola pasada y reutilizando la memoria entre cuadros (Sección~\ref{sec:res-rendimiento}).

\subsection{Detección de manos con MediaPipe}

El detector se configura para hasta dos manos, en modo de imagen (cada cuadro se procesa de forma independiente), con una confianza mínima de detección configurable que por defecto es 0,6. Se ejecuta en la GPU cuando el dispositivo lo permite y, si la inicialización falla, la aplicación vuelve automáticamente al procesador central. El cuadro llega al detector sin espejar; el espejado de la cámara frontal se aplica solo a la vista previa y al dibujo de los puntos en pantalla.

\subsection{Construcción de características}

Por cada cuadro con manos se arma un vector de 126 valores: las coordenadas de los 21 puntos de la mano izquierda seguidas de las de la mano derecha, con ceros en la mano que no se detecta y la corrección de las coordenadas a un lienzo cuadrado (Sección~\ref{sec:representacion}). Los vectores se acumulan durante la captura y, al completarla, la secuencia se lleva a 15 cuadros.

El contrato de características está implementado en Kotlin y verificado con los mismos casos de prueba que la implementación de Python. Lo usan el entrenamiento en el dispositivo y el esquema federado; la traducción en vivo, en cambio, todavía arma la secuencia con el procedimiento anterior, lo que constituye una de las limitaciones del trabajo (Capítulo~\ref{cap:discusion}).

\subsection{Clasificación con TensorFlow Lite}

El clasificador carga el modelo y la lista de clases (\texttt{words.json}), cuyo orden coincide con el de las salidas del modelo. Si el dispositivo recibió un modelo actualizado por el esquema federado, se usa ese; si no, o si el paquete recibido está dañado, se usa el modelo de fábrica incluido en la aplicación.

La salida del modelo se ajusta con la temperatura de calibración ($T = 1{,}20$), y una predicción se muestra solo si su confianza calibrada supera el umbral, que por defecto es 0,75 y puede modificarse en los ajustes (Sección~\ref{sec:res-calibracion}). La clase «nada» nunca se muestra ni se pronuncia. Cuando una captura termina en «nada» o por debajo del umbral, la pantalla indica que la seña no se reconoció y sugiere repetirla.

\subsection{Salida en texto y voz}

Cada seña reconocida se agrega a una frase de hasta seis palabras que se muestra en la parte inferior de la pantalla; una misma seña repetida en dos capturas seguidas no se duplica. La palabra se pronuncia con el motor de síntesis de Android, en la variante de español elegida (por defecto, la de Argentina) y con la velocidad y el tono configurados. La voz dice exactamente lo que muestra la pantalla: las predicciones descartadas no se pronuncian. La frase puede repetirse en voz alta o borrarse con un botón.

\section{Modo voz a texto y escenario con dos dispositivos}

El modo voz a texto transcribe lo que dice la persona oyente para que la persona sorda lo lea. Usa el reconocedor de voz del sistema con resultados parciales, de modo que el texto aparece mientras la persona habla, y le pide resolver el reconocimiento en el dispositivo. Esa indicación es una preferencia que el servicio puede no respetar; por eso, cuando el modo local está activo, la aplicación no reintenta en la nube si el reconocimiento falla, sino que informa que falta el paquete de voz del idioma. Entre el reconocimiento y la pantalla se aplica la limpieza del texto descrita en el Capítulo~\ref{cap:desafios}.

El escenario de uso previsto combina dos dispositivos con la misma aplicación en modos complementarios (Figura~\ref{fig:dos-dispositivos}): la persona sorda señala frente a su teléfono, que muestra y pronuncia la traducción, y la persona oyente responde hablando frente al suyo, que transcribe lo que dice para que la persona sorda lo lea.

\begin{figure}[h]
    \centering
    \begin{tikzpicture}[
        persona/.style={draw, circle, minimum size=1.3cm, align=center, font=\small},
        disp/.style={draw, rounded corners, minimum width=2.8cm, minimum height=1.2cm, align=center, font=\small},
        flecha/.style={-{Stealth}, thick}
    ]
        \node[persona] (sorda) {Persona\\sorda};
        \node[disp, right=1.2cm of sorda, yshift=0.9cm] (d1) {Teléfono A\\Traducir señas};
        \node[disp, right=1.2cm of sorda, yshift=-0.9cm] (d2) {Teléfono B\\Traducir voz};
        \node[persona, right=1.2cm of d1, yshift=-0.9cm] (oyente) {Persona\\oyente};
        \draw[flecha] (sorda) -- node[above, sloped, font=\scriptsize] {señas} (d1.west);
        \draw[flecha] (d1.east) -- node[above, sloped, font=\scriptsize] {texto y voz} (oyente);
        \draw[flecha] (oyente) -- node[below, sloped, font=\scriptsize] {habla} (d2.east);
        \draw[flecha] (d2.west) -- node[below, sloped, font=\scriptsize] {texto} (sorda);
    \end{tikzpicture}
    \caption{Escenario de uso con dos dispositivos en modos complementarios.}
    \label{fig:dos-dispositivos}
\end{figure}

\section{Modo de entrenamiento y recolección de muestras}

El modo de entrenamiento permite grabar tomas etiquetadas desde la propia aplicación. La persona elige una seña de la lista, que muestra cuántas tomas tiene cada una, y la aplicación graba sola cada vez que detecta una seña completa, con la misma máquina de estados de la traducción. Una imagen de referencia de la seña queda visible en una esquina mientras se graba, un contador indica el avance y un botón permite deshacer la última toma. También es posible agregar una seña nueva, que no forma parte del modelo, escribiendo su nombre.

En la sesión guiada (Sección~\ref{sec:protocolo-captura}), la aplicación lee la planilla con el orden de las tomas de cada participante y avanza sola de una seña a la siguiente, con pausas entre tandas y una indicación de cómo debe quedar cada condición. Si la aplicación se cierra a mitad de la sesión, se retoma donde quedó.

Cada toma se guarda en el almacenamiento propio de la aplicación en dos formatos: el formato original, con la secuencia ya procesada, y el formato v2 del conjunto de datos, que conserva los landmarks crudos de cada cuadro y sus metadatos (Sección~\ref{sec:representacion}). Las tomas se asocian a un seudónimo derivado de la cuenta, que no permite recuperar la identidad de la persona. Una pantalla de gestión del conjunto de datos muestra las tomas por seña y permite exportarlas mediante el menú de compartir del sistema, para incorporarlas al reentrenamiento.

\section{Aprendizaje federado}

La pantalla de aprendizaje federado reúne las funciones descritas en el Capítulo~\ref{cap:federado}: buscar una versión nueva del modelo, participar en una ronda de entrenamiento, aportar las propias tomas al servidor, personalizar el modelo con las tomas propias y volver al modelo de fábrica. Antes de participar o aportar, la aplicación explica en lenguaje simple qué datos salen del teléfono y pide confirmación.

\section{Interfaz de usuario}

La pantalla principal ofrece las dos funciones de traducción, señas y voz, y debajo, el diccionario de señas y el modo de entrenamiento. Las pantallas son:

\begin{itemize}
    \item \textbf{Traducir señas.} La vista de la cámara ocupa toda la pantalla, con un recuadro que indica el encuadre y se resalta cuando se detectan manos, los puntos de las manos dibujados encima y, en la parte inferior, la frase traducida y un rótulo de estado («esperando señas», «leyendo la seña», «no la reconocí»).
    \item \textbf{Traducir voz.} La transcripción en curso y las frases anteriores, en letra grande.
    \item \textbf{Diccionario.} Las señas que reconoce el modelo, organizadas en palabras y letras, con un buscador y una imagen de referencia de cada seña; «mi nombre es» y «por favor» todavía no tienen imagen.
    \item \textbf{Modo entrenamiento y conjunto de datos.} La grabación de tomas y su gestión.
    \item \textbf{Aprendizaje federado.} Las funciones de la sección anterior.
    \item \textbf{Ajustes.} Resolución de análisis, cámara, sensibilidad de detección, dibujo de los puntos, lectura en voz alta, velocidad y tono de la voz, variante del español, umbral de confianza y reconocimiento de voz solo local.
\end{itemize}

Las Figuras~\ref{fig:pantallas-principales} a~\ref{fig:pantallas-ajustes} muestran las pantallas principales, capturadas en la tableta Lenovo con la cámara apuntando al techo, para que no aparezca ninguna persona.

\newcommand{\captura}[2]{%
    \begin{minipage}[t]{0.31\textwidth}
        \centering
        \includegraphics[width=\linewidth]{capturas/#1}\\[2pt]
        {\small #2}
    \end{minipage}%
}

\begin{figure}[p]
    \centering
    \captura{bienvenida}{(a) Bienvenida}\hfill
    \captura{inicio}{(b) Inicio}\hfill
    \captura{senas}{(c) Traducir señas}
    \caption{Pantalla de bienvenida, pantalla principal y traducción de señas en espera, con el recuadro de encuadre y el rótulo de estado.}
    \label{fig:pantallas-principales}
\end{figure}

\begin{figure}[p]
    \centering
    \captura{voz}{(a) Traducir voz}\hfill
    \captura{diccionario}{(b) Diccionario}\hfill
    \captura{federado}{(c) Aprendizaje federado}
    \caption{Modo voz a texto, diccionario de señas y pantalla de aprendizaje federado, con el texto que explica qué datos salen del dispositivo.}
    \label{fig:pantallas-voz}
\end{figure}

\begin{figure}[p]
    \centering
    \captura{entrenamiento_inicio}{(a) Modo entrenamiento}\hfill
    \captura{entrenamiento_lista}{(b) Elección de la seña}\hfill
    \captura{entrenamiento_camara}{(c) Grabación}
    \caption{Modo entrenamiento: menú, lista de señas con la cantidad de tomas de cada una y pantalla de grabación, con la imagen de referencia y el contador.}
    \label{fig:pantallas-entrenamiento}
\end{figure}

\begin{figure}[p]
    \centering
    \captura{dataset}{(a) Conjunto de datos}\hspace{0.04\textwidth}
    \captura{ajustes}{(b) Ajustes}
    \caption{Gestión del conjunto de datos, con la exportación, y parte de la pantalla de ajustes.}
    \label{fig:pantallas-ajustes}
\end{figure}

\pendienteautora{Imágenes de referencia del diccionario: reemplazarlas por imágenes propias o citar su fuente.}

El núcleo comprometido en la propuesta es el pipeline de traducción de señas, la salida en texto y voz y el modo voz a texto. Durante el desarrollo se incorporaron mejoras de experiencia de uso que surgieron de las pruebas: la presentación inicial, la precarga de los modelos mientras se muestra la pantalla de inicio para que la primera traducción no espere, el rótulo de estado de la traducción, la pantalla que permanece encendida durante la traducción y la grabación, y la pantalla de cámara fija en orientación vertical (Capítulo~\ref{cap:desafios}).

\chapter{Aprendizaje federado}
\label{cap:federado}

Este capítulo describe el esquema de aprendizaje federado incorporado al prototipo: su arquitectura, el protocolo de comunicación, el entrenamiento en el dispositivo y las medidas de protección de las actualizaciones. Los resultados se presentan en la Sección~\ref{sec:res-federado}.

\section{Motivación y alcance}

La escasez de datos anotados de LSA limita la capacidad del modelo de reconocer a personas que no participaron de su entrenamiento: el acierto cae de 84,6\,\% con una persona conocida a 67,3\,\% con una persona no vista (Sección~\ref{sec:res-escenarios}). Mejorar el modelo requiere incorporar tomas de nuevas personas, pero centralizar grabaciones de gestos implica recolectar datos personales.

El aprendizaje federado permite entrenar un modelo de forma colaborativa a partir de actualizaciones calculadas en cada dispositivo, sin transferir los datos crudos~\cite{mcmahan2017fedavg}. En este trabajo se aborda con carácter experimental: se implementa el ciclo completo en dispositivos reales y se evalúa su viabilidad, sin pretender un mecanismo de producción.

\section{Arquitectura}

El esquema se compone de un servidor y de la aplicación en cada dispositivo. Para la tesina, el servidor se ejecuta en una computadora personal dentro de una red local, sin acceso a internet. El mismo código puede trasladarse a un servidor en la nube; esa migración queda como trabajo futuro.

El modelo se divide en dos partes:

\begin{itemize}
    \item La parte de extracción de características (dos capas LSTM), que transforma la secuencia de $15 \times 126$ valores en un vector de 128 valores. Se distribuye como modelo TensorFlow Lite y no se entrena en el dispositivo.
    \item La cabeza de clasificación (tres capas densas de 64, 64 y $C$ neuronas, con $C = 13$), de 13\,261 parámetros. Es la parte que se entrena en el dispositivo y que viaja en las rondas.
\end{itemize}

La aplicación clasifica las señas con la parte de extracción y la cabeza vigente. La cabeza se implementó en Kotlin, con propagación hacia atrás y el optimizador Adam, porque entrenar en el dispositivo con TensorFlow Lite exige operaciones que el núcleo de la biblioteca no incluye. La implementación se verificó contra TensorFlow con una tolerancia de $10^{-5}$ en la salida, la pérdida, los gradientes y los pesos tras uno y tres pasos de optimización.

\pendiente{Agregar un diagrama de la arquitectura: servidor en la red local, dispositivos, rondas y paquetes.}

\section{Ciclo de una ronda}

Cada ronda $r$ sigue estos pasos:

\begin{enumerate}
    \item Cada dispositivo descarga la cabeza vigente $w_r$ y verifica su integridad mediante un resumen SHA-256.
    \item Entrena la cabeza durante dos épocas con sus propias tomas y calcula la actualización $\Delta_k = w_k - w_r$.
    \item Cuantiza $\Delta_k$ a 8 bits y la envía al servidor junto con la cantidad de tomas $n_k$ utilizadas. Las tomas no salen del dispositivo.
    \item Cuando se reúne la cantidad mínima de participantes, o vence el plazo de la ronda, el servidor combina las actualizaciones mediante el promedio ponderado de FedAvg:
    \begin{equation}
        w_{r+1} = w_r + \sum_{k} \frac{n_k}{n}\,\Delta_k, \qquad n = \sum_k n_k.
    \end{equation}
    \item El servidor evalúa $w_{r+1}$ sobre un conjunto de prueba fijo y la publica solo si su acierto no baja más de un punto respecto de $w_r$.
    \item Cada dispositivo descarga la nueva cabeza al buscar una actualización.
\end{enumerate}

El conjunto de prueba fijo está formado por el 15\,\% de las tomas de los dispositivos, estratificado por seña. Esas tomas no se utilizan para entrenar ni en el servidor ni en los dispositivos: el servidor publica sus identificadores y la aplicación las excluye.

\section{Compresión y protección de las actualizaciones}

\subsection{Cuantización a 8 bits}

Cada tensor de la actualización se cuantiza de forma simétrica con su propia escala $s = \max|x| / 127$. El redondeo es estocástico: el valor $x/s$ se redondea hacia arriba con probabilidad igual a su parte fraccionaria, lo que hace que la cuantización sea insesgada. La actualización ocupa 13\,285 bytes frente a 53\,044 en punto flotante de 32 bits.

Para garantizar que el servidor reconstruya exactamente lo que el dispositivo envió, la cuantización de Kotlin y la de Python se verifican bit a bit sobre un conjunto de casos compartido entre ambos repositorios.

\subsection{Límite a la magnitud}

La propuesta plantea limitar la magnitud de las actualizaciones para evitar que un dispositivo con datos erróneos degrade el modelo de todos. El servidor escala cualquier actualización cuya norma L2 supere un límite $L$:
\begin{equation}
    \Delta_k \leftarrow \Delta_k \cdot \min\left(1, \frac{L}{\lVert \Delta_k \rVert_2}\right).
\end{equation}
Se adoptó $L = 2{,}0$, tres veces la mediana de las actualizaciones honestas medidas con las tres personas.

\subsection{Rotación aleatoria}

Siguiendo las actualizaciones esbozadas de Kone\v{c}n\'y et al.~\cite{konecny2016}, se implementó una rotación aleatoria previa a la cuantización: $y = H D x / \sqrt{n}$, donde $D$ es una matriz diagonal de signos aleatorios generados a partir de una semilla que viaja con la actualización y $H$ es la transformada de Walsh-Hadamard. La rotación reparte la energía del vector entre todos sus valores, de modo que un valor extremo no fije la escala de todo el tensor. Como con las actualizaciones de este modelo no redujo el error (Sección~\ref{sec:res-federado}), se mantiene implementada pero desactivada.

\section{Reentrenamiento completo y distribución de modelos}

Además de las rondas sobre la cabeza, el servidor puede reentrenar el modelo completo. Las personas pueden aportar sus tomas de forma opcional, con consentimiento explícito: se envían solo las coordenadas de los puntos de las manos, nunca imágenes, asociadas a un identificador seudónimo derivado de la cuenta.

El servidor reentrena con todas las fuentes, verifica que el modelo nuevo no tenga operaciones Flex, que la parte de extracción en TensorFlow Lite coincida con la de Keras y que el acierto sobre el conjunto de prueba fijo no empeore, y lo publica como un paquete versionado. El paquete incluye la parte de extracción, la cabeza, la calibración y un manifiesto con el resumen SHA-256 de cada archivo. La aplicación descarga el paquete, verifica cada archivo antes de usarlo y conserva siempre el modelo de fábrica para poder volver a él.

Si las tomas aportadas incluyen una seña nueva, el paquete sale con una clase más. Esto permite ampliar el vocabulario sin publicar una nueva versión de la aplicación.

\section{Personalización en el dispositivo}

FedAvg produce un modelo para el conjunto de los participantes. Para adaptarlo a cada persona, la aplicación permite personalizar la cabeza: parte de la cabeza compartida vigente, la entrena durante ocho épocas con las tomas propias y la guarda solo en el dispositivo. La aplicación traduce con la cabeza personalizada mientras corresponda a la versión compartida vigente; cuando llega una versión nueva, se vuelve a personalizar automáticamente. Las rondas federadas parten siempre de la cabeza compartida.

\section{Privacidad}

En las rondas solo sale del dispositivo la actualización cuantizada de la cabeza y la cantidad de tomas utilizadas. El aporte de tomas para el reentrenamiento completo es una opción separada que requiere consentimiento explícito y envía solo las coordenadas de los puntos de las manos. La traducción de señas no utiliza la red en ningún caso.

En la versión de depuración utilizada para el piloto, la comunicación con el servidor se realiza por HTTP sin cifrar dentro de la red local; la versión de distribución bloquea ese tráfico. La comunicación cifrada, la agregación segura y la privacidad diferencial quedan como trabajo futuro.

\section{Decisiones de implementación}

La propuesta mencionaba TensorFlow Federated y Flower como marcos posibles. Se optó por una implementación propia, de unas pocas centenas de líneas, por dos motivos: TensorFlow Federated no ofrece versiones para Windows y es incompatible con la versión de TensorFlow utilizada, y el cliente de Flower para Android impone su propio formato de modelo, que no se ajusta al entrenamiento de la cabeza en Kotlin. La implementación propia se verificó contra una simulación de FedAvg en Python.

\chapter{Desafíos técnicos y decisiones de diseño}
\label{cap:desafios}

% Cada sección sigue el esquema: problema observado, causa identificada,
% alternativas evaluadas, solución adoptada e impacto.

\section{Espejado de las manos entre la cámara, MediaPipe y el modelo}

\pendiente{Describir la inconsistencia entre la imagen espejada de la cámara frontal, la lateralidad (handedness) informada por MediaPipe y la disposición de los datos con que se entrenó el modelo; cómo se detectó, cómo se resolvió y qué efecto tuvo en el reconocimiento.}

\section{Rendimiento en dispositivos de gama baja}

\pendiente{Describir las dificultades observadas al ejecutar la aplicación en dispositivos antiguos, la identificación del cuello de botella (captura, conversión de imagen, detección de landmarks o inferencia) y las optimizaciones aplicadas.}

\section{Orientación de la cámara y estabilidad del seguimiento}

\pendiente{Describir el efecto de la orientación del dispositivo sobre la estabilidad de los landmarks.}

\section{Otros desafíos}

\pendiente{Agregar otros problemas relevantes encontrados durante el desarrollo, por ejemplo la limpieza del texto reconocido o los cambios de nomenclatura de las señas.}

\chapter{Experimentos y resultados}

\section{Desempeño del modelo}

\pendiente{Reportar accuracy, F1 macro y micro y matriz de confusión sobre el conjunto de prueba. Completar cuando el entrenamiento esté cerrado.}

\section{Consistencia estadística}

\pendiente{Reportar la métrica de consistencia $\sigma$ definida en la propuesta.}

\section{Robustez ante condiciones variables}

\pendiente{Reportar resultados según usuario, iluminación, ángulo y distancia de cámara.}

\section{Rendimiento en distintos dispositivos}

\pendiente{Completar la tabla con los seis dispositivos probados (tres antiguos, un teléfono reciente y dos tablets de alta gama), indicando condiciones de medición. Los valores actuales son preliminares y deben revalidarse con el modelo final.}

\begin{table}[h]
    \centering
    \begin{tabular}{llll}
        \toprule
        Dispositivo & SoC & FPS efectivos & Inferencia \\
        \midrule
        Galaxy Tab S7 & Snapdragon 865+ & 14.0 & 32 ms \\
        Galaxy A55 & Exynos 1480 & 5.3 & 30 ms \\
        Galaxy A31 & Helio P65 & 3.0 & 559 ms \\
        \bottomrule
    \end{tabular}
    \caption{Resultados preliminares de rendimiento en distintos dispositivos.}
    \label{tab:rendimiento-dispositivos}
\end{table}

\section{Comparación entre variantes del modelo}

\pendiente{Comparar tamaño, latencia y precisión entre la variante float32 y la variante optimizada.}

\chapter{Discusión}

\section{Interpretación de los resultados}

\pendiente{Analizar qué muestran los resultados del modelo y del rendimiento, sin repetirlos.}

\section{Requisitos mínimos de hardware}

\pendiente{A partir de las pruebas en los seis dispositivos, establecer qué características de hardware permiten un uso interactivo aceptable y cuáles no.}

\section{Comparación con trabajos relacionados}

\pendiente{Contrastar los resultados con los antecedentes revisados en el Capítulo 2, explicitando las diferencias de condiciones.}

\section{Limitaciones y amenazas a la validez}

\pendiente{Tamaño y diversidad del conjunto de datos, cantidad de personas que realizaron las señas, vocabulario acotado y condiciones de prueba.}

\chapter{Conclusiones y trabajo futuro}

\section{Conclusiones}

\pendiente{Retomar el objetivo general y cada objetivo específico, indicando en qué medida se cumplieron.}

\section{Aportes del trabajo}

\begin{itemize}
    \item Desarrollo de un prototipo funcional de traducción de LSA a texto y voz.
    \item Integración de visión por computadora y aprendizaje automático en Android.
    \item Ejecución local orientada a privacidad y disponibilidad.
    \item Evaluación en distintos dispositivos móviles.
\end{itemize}

\section{Trabajo futuro}

\begin{itemize}
    \item Ampliar el conjunto de señas reconocidas.
    \item Mejorar y diversificar el conjunto de datos.
    \item Optimizar la conversión de imagen y el procesamiento de landmarks.
    \item Evaluar modelos más robustos para secuencias temporales.
\end{itemize}


\clearpage
\addcontentsline{toc}{chapter}{Bibliografía}
\bibliographystyle{unsrt}
\bibliography{bibliografia}

\appendix
\chapter{Anexos}
En esta sección pueden incorporarse capturas de pantalla, tablas complementarias, detalles del conjunto de datos, etiquetas del modelo y fragmentos relevantes de configuración.

\end{document}
